\documentclass[12pt, a4paper]{ctexart}
\usepackage{amsmath, amssymb, amsthm}
\usepackage{geometry}
\usepackage{fancyhdr}
\usepackage{enumitem}
\usepackage{tikz}
\usepackage{pgfplots}
\usepackage{array}
\usepackage{multirow}
\usepackage{booktabs}

\geometry{left=2cm, right=2cm, top=2.5cm, bottom=2.5cm}

\setCJKmainfont{Hiragino Sans GB}
\setCJKsansfont{Heiti SC}

\pagestyle{fancy}
\fancyhf{}
\fancyhead[L]{直线和圆的方程 测试卷}
\fancyhead[R]{数学（选择性必修第一册）}
\fancyfoot[C]{第 \thepage\ 页}
\renewcommand{\headrulewidth}{0.4pt}

\setlength{\parindent}{0pt}

\newcommand{\score}[1]{\hfill\textbf{（#1分）}}

\begin{document}

\begin{center}
{\Large\bfseries 直线和圆的方程 测试卷}\\[0.5em]
{\normalsize 时间：120分钟 \quad 满分：150分}
\end{center}

\vspace{1em}

\noindent\textbf{一、单项选择题（本大题共8小题，每小题5分，共40分。每小题只有一个选项符合题目要求。）}

\vspace{0.5em}

\begin{enumerate}
\item 过点 $(1,2)$ 且斜率为 $3$ 的直线方程为 \score{5}
\begin{enumerate}[label=\Alph*.]
\item $y=3x-1$
\item $y=3x+1$
\item $y=3x-2$
\item $y=3x+2$
\end{enumerate}

\item 直线 $2x-3y+6=0$ 在 $x$ 轴上的截距为 \score{5}
\begin{enumerate}[label=\Alph*.]
\item $-3$
\item $3$
\item $-2$
\item $2$
\end{enumerate}

\item 两条直线 $l_1: x+2y-1=0$ 与 $l_2: 2x-y+3=0$ 的位置关系是 \score{5}
\begin{enumerate}[label=\Alph*.]
\item 平行
\item 垂直
\item 相交但不垂直
\item 重合
\end{enumerate}

\item 过点 $(1,-1)$ 且与直线 $3x-y+2=0$ 平行的直线方程为 \score{5}
\begin{enumerate}[label=\Alph*.]
\item $3x-y-4=0$
\item $3x-y-2=0$
\item $x+3y+2=0$
\item $3x+y-2=0$
\end{enumerate}

\item 点 $(3,4)$ 到直线 $3x-4y+5=0$ 的距离为 \score{5}
\begin{enumerate}[label=\Alph*.]
\item $2$
\item $\dfrac{4}{5}$
\item $4$
\item $\dfrac{8}{5}$
\end{enumerate}

\item 圆 $x^2+y^2-4x+6y-3=0$ 的圆心和半径分别为 \score{5}
\begin{enumerate}[label=\Alph*.]
\item $(2,-3)$，$4$
\item $(-2,3)$，$4$
\item $(2,-3)$，$16$
\item $(-2,3)$，$2$
\end{enumerate}

\item 过点 $A(1,0)$，$B(0,1)$，$C(2,1)$ 的圆的方程为 \score{5}
\begin{enumerate}[label=\Alph*.]
\item $(x-1)^2+(y-1)^2=1$
\item $(x-1)^2+y^2=1$
\item $x^2+(y-1)^2=1$
\item $(x-\dfrac{3}{2})^2+(y-\dfrac{1}{2})^2=\dfrac{1}{2}$
\end{enumerate}

\item 直线 $y=x+b$ 与圆 $x^2+y^2=4$ 相切，则 $b$ 的值为 \score{5}
\begin{enumerate}[label=\Alph*.]
\item $\pm\sqrt{2}$
\item $\pm 2$
\item $\pm 2\sqrt{2}$
\item $\pm 4$
\end{enumerate}
\end{enumerate}

\vspace{1em}

\noindent\textbf{二、多项选择题（本大题共4小题，每小题5分，共20分。每小题有多个选项符合题目要求，全部选对得5分，部分选对得2分，有选错的得0分。）}

\vspace{0.5em}

\begin{enumerate}[resume]
\item 关于直线 $l: ax+by+c=0$（$a$，$b$ 不同时为零），下列结论正确的是 \score{5}
\begin{enumerate}[label=\Alph*.]
\item 当 $a=0$ 时，直线平行于 $x$ 轴
\item 当 $b=0$ 时，直线平行于 $y$ 轴
\item 直线的方向向量为 $(b,-a)$
\item 直线的斜率为 $-\dfrac{a}{b}$（$b\neq 0$ 时）
\end{enumerate}

\item 已知直线 $l_1: (a+2)x+(1-a)y-3=0$ 与 $l_2: (a-1)x+(2a+3)y+2=0$ 互相垂直，则 $a$ 的值可能为 \score{5}
\begin{enumerate}[label=\Alph*.]
\item $1$
\item $-1$
\item $-2$
\item $\dfrac{1}{2}$
\end{enumerate}

\item 关于圆 $(x-1)^2+(y+2)^2=9$，下列结论正确的是 \score{5}
\begin{enumerate}[label=\Alph*.]
\item 圆心为 $(1,-2)$
\item 半径为 $3$
\item 点 $(1,1)$ 在圆内
\item 圆经过原点
\end{enumerate}

\item 已知圆 $C_1: x^2+y^2=4$ 与圆 $C_2: (x-3)^2+(y-4)^2=r^2$，下列结论正确的是 \score{5}
\begin{enumerate}[label=\Alph*.]
\item 当 $r=3$ 时，两圆相交
\item 当 $r=5$ 时，两圆外切
\item 当 $r=7$ 时，两圆内含
\item 当 $r=2$ 时，两圆内切
\end{enumerate}
\end{enumerate}

\vspace{1em}

\noindent\textbf{三、填空题（本大题共4小题，每小题5分，共20分。）}

\vspace{0.5em}

\begin{enumerate}[resume]
\item 经过两点 $A(2,3)$，$B(5,-1)$ 的直线的斜率 $k=$ \underline{\hspace{2cm}}，直线方程为 \underline{\hspace{4cm}}。\score{5}

\item 若直线 $l$ 过点 $(2,-3)$，且在两坐标轴上的截距相等，则直线 $l$ 的方程为 \underline{\hspace{5cm}}。\score{5}

\item 圆 $x^2+y^2-6x+8y=0$ 关于直线 $y=x$ 对称的圆的方程为 \underline{\hspace{5cm}}。\score{5}

\item 过点 $(1,2)$ 向圆 $(x-3)^2+(y-4)^2=1$ 引切线，切线长为 \underline{\hspace{2cm}}。\score{5}
\end{enumerate}

\vspace{1em}

\noindent\textbf{四、解答题（本大题共6小题，共70分。解答应写出文字说明、证明过程或演算步骤。）}

\vspace{0.5em}

\begin{enumerate}[resume]
\item （10分）已知三角形的三个顶点为 $A(1,2)$，$B(3,4)$，$C(-1,6)$。
\begin{enumerate}[label=(\arabic*)]
\item 求 $BC$ 边所在直线的方程；
\item 求 $BC$ 边上的高所在直线的方程；
\item 求三角形 $ABC$ 的面积。
\end{enumerate}

\item （10分）已知直线 $l_1: 2x+y-4=0$，$l_2: x-2y+3=0$。
\begin{enumerate}[label=(\arabic*)]
\item 求 $l_1$ 与 $l_2$ 的交点坐标；
\item 求过 $l_1$ 与 $l_2$ 的交点且与直线 $3x-y+1=0$ 平行的直线方程；
\item 求 $l_1$ 关于 $l_2$ 对称的直线方程。
\end{enumerate}

\item （12分）已知圆 $C$ 经过点 $A(1,3)$，$B(5,1)$，且圆心在直线 $y=x$ 上。
\begin{enumerate}[label=(\arabic*)]
\item 求圆 $C$ 的方程；
\item 判断点 $M(2,4)$ 与圆 $C$ 的位置关系；
\item 求过点 $M(2,4)$ 的最短弦所在的直线方程。
\end{enumerate}

\item （12分）已知圆 $O: x^2+y^2=9$ 和直线 $l: x-y+4=0$。
\begin{enumerate}[label=(\arabic*)]
\item 求圆心到直线 $l$ 的距离；
\item 求直线 $l$ 被圆 $O$ 截得的弦长；
\item 求与直线 $l$ 平行且与圆 $O$ 相切的直线方程。
\end{enumerate}

\item （12分）已知点 $P$ 在圆 $x^2+y^2=4$ 上运动，点 $Q(4,0)$。
\begin{enumerate}[label=(\arabic*)]
\item 求 $|PQ|$ 的最大值和最小值；
\item 求线段 $PQ$ 中点 $M$ 的轨迹方程；
\item 求 $\dfrac{y}{x-4}$（$P$ 点坐标为 $(x,y)$）的取值范围。
\end{enumerate}

\item （14分）已知圆 $C: x^2+y^2-4x-2y-20=0$，直线 $l: (2m+1)x+(m+1)y-7m-4=0$（$m\in\mathbb{R}$）。
\begin{enumerate}[label=(\arabic*)]
\item 证明：直线 $l$ 恒过定点，并求出该定点坐标；
\item 判断定点与圆 $C$ 的位置关系；
\item 当直线 $l$ 被圆 $C$ 截得的弦长最短时，求 $m$ 的值及最短弦长；
\item 设圆 $C$ 与 $x$ 轴交于 $A$，$B$ 两点，求以 $AB$ 为直径的圆的方程。
\end{enumerate}
\end{enumerate}

\end{document}
